\section{Related Work}

\subsection{Pre-training Corpus Curation and Model Generalization}

The construction of large-scale pre-training corpora has evolved substantially, with increasing emphasis on quality filtering, deduplication, and domain balancing. The Pile~\cite{Gao2020Pile} established a multi-domain composition strategy, mixing web text with curated sources such as PubMed, GitHub, and legal text. Pythia~\cite{Biderman2023Pythia} built on The Pile to provide publicly released training checkpoints at multiple scales, enabling training-dynamics research. Dolma~\cite{Soldaini2024Dolma} advanced corpus curation further, applying multi-stage deduplication, quality classifiers, and explicit inclusion of academic sources (S2ORC, Wikipedia) under the hypothesis that these choices yield more generalizable models. RefinedWeb~\cite{Penedo2023RefinedWeb} demonstrated that aggressive deduplication of web-only data can match or exceed mixed-source corpora on downstream benchmarks. OLMo~\cite{Groeneveld2024OLMo} adopted Dolma as its training corpus, releasing full training details and intermediate checkpoints to support reproducible research. Scaling studies~\cite{Muennighoff2023Scaling} have examined how data quantity and quality interact at different compute budgets, but matched-architecture, matched-scale comparisons of corpus quality in isolation remain rare.

\subsection{Generalization Evaluation Metrics}

MMLU~\cite{Hendrycks2021MMLU} measures factual knowledge across 57 academic subjects and is widely used as a proxy for knowledge acquisition from pre-training data. HellaSwag~\cite{Zellers2019HellaSwag} evaluates commonsense reasoning via sentence completion and serves as a complementary generalization signal. ARC~\cite{Clark2018ARC} provides science question answering at two difficulty levels (Easy and Challenge), probing reasoning beyond surface pattern matching. The use of ratio metrics combining these benchmarks to measure generalization balance has been explored in the context of understanding knowledge vs.\ commonsense trade-offs~\cite{Miller2021Accuracy,Koh2021WILDS}. However, the sensitivity of such ratio metrics to denominator saturation --- where one benchmark component ceases to vary meaningfully across model variants --- has not been systematically characterized.

\subsection{Data Attribution and Quality Proxy Methods}

A parallel line of research addresses the question of which training examples are responsible for model behavior, using influence functions~\cite{Koh2017Influence}, TRAK~\cite{Park2023TRAK}, and TracIn~\cite{Pruthi2020TracIn}. These methods provide fine-grained attribution but are computationally expensive at pre-training scale. Membership inference methods~\cite{shi2024detecting} have been applied to detect training data contamination in benchmark evaluations --- a concern directly relevant to MMLU, which overlaps with web-crawled corpora. Our work operates at a coarser granularity, using aggregate benchmark performance as a quality proxy, and identifies structural limitations of this approach.

\subsection{Positioning Our Contribution}

Our work differs from prior corpus comparison studies in three respects. First, we fix training scale ($\sim$300B tokens) using documented public checkpoints rather than training new models, enabling reproducibility without compute cost. Second, we use a ratio metric specifically designed to capture the balance between knowledge acquisition (MMLU) and commonsense generalization (HellaSwag), rather than reporting benchmark scores in isolation. Third, we report a null and negative result --- the curated-corpus model does not achieve a higher generalization balance ratio --- and analyze the structural reason: HellaSwag denominator saturation that limits the metric's discriminative power at this scale. This analysis motivates concrete methodological recommendations for future corpus quality studies.
